% §6.7 Cross-method orthogonality between Shen, Baur, and counterfactual.
% Numbers locked against results/replications/cross_method_concordance.csv.
% Spearman and sign-agreement computed on the per-city point estimates.

\subsection{Cross-method orthogonality}
\label{sec:cross-method-12}

Sections~\ref{sec:shen-12city} and \ref{sec:baur-12city} each identify
a different small subset of metropolitan markets in which a
text-on-price effect can be detected at conventional significance.
The natural question is whether the two channels agree on which
markets have an effect, or whether they identify substantively
different aspects of listing language. We answer this question by
computing per-market rank and sign correlations between the
Shen-Ross uniqueness DML estimates and the Baur PC1 DML estimates,
and we add the counterfactual rewriting total effect from
Section~\ref{sec:counterfactual-12} as a third axis.

\paragraph{Pairwise Spearman correlations.}
Table~\ref{tab:cross-method-12} reports the per-market point
estimates for the three channels under the pooled-PCA Baur
specification of Section~\ref{sec:baur-12city}. Across the twelve
markets in our panel, the Spearman rank correlation between the
Shen-Ross point estimate and the Baur point estimate is $-0.15$
($p = 0.65$), indistinguishable from zero. The Spearman correlation
between the Shen-Ross estimate and the counterfactual total effect
is $-0.07$ ($p = 0.83$), also indistinguishable from zero. The
correlation between the Baur estimate and the counterfactual total
effect is $+0.48$ ($p = 0.12$), suggestive but not statistically
significant at conventional levels.  The Baur-counterfactual
positive correlation is consistent with the substantive reading that
both channels measure aspects of language-induced price change in
the substantively consistent direction (utilitarian or starter-home
language correlating with lower prices), while the Shen-Ross
uniqueness channel captures a logically distinct property of
listing language (distinctiveness from local peer set) whose
correlation with the other two is small in absolute magnitude.

\paragraph{Sign agreement.} The three methods agree on a single
common sign of the per-market point estimate in $3$ of the $12$
markets. Pairwise, Shen and Baur agree on sign in $4$ of $12$
markets, Shen and the counterfactual in $5$ of $12$, and Baur and
the counterfactual in $5$ of $12$.  The pooled-PCA Baur axis returns
negative point estimates in eleven of the twelve markets (Atlanta
the sole exception with a near-zero positive estimate), so most
sign-disagreements with Shen reflect the difference between
distinctive-language premia (Shen) and utilitarian-axis discounts
(Baur), rather than methodological inconsistency.

\paragraph{Substantive interpretation under a triangulation
design.} The three identification strategies are not redundant point
estimators of a single scalar effect; they are deliberately chosen
to have orthogonal sources of bias.  The Shen-Ross uniqueness measure
leverages within-market lexical rarity and inherits any bias from the
neighborhood-peer-set definition.  The pooled-PCA Baur sentence-BERT
axis leverages cross-listing semantic geometry and inherits bias from
the embedding model's pretraining distribution.  The LLM counter-
factual rewriting intervention leverages within-listing minimal-edit
substitution and inherits bias from the rewriter's distributional
shift.  Following the causal-triangulation framework of
\citet{lawlor2016triangulation} and \citet{hammerton2021causal},
methods whose biases are largely unrelated provide jointly stronger
evidence on a common causal target than methods with shared bias
structure, regardless of whether their per-market point estimates
agree numerically.  This is a formal restatement of the Heckman
\citep{heckman2008econometric} position that "different estimators
identify different parameters," reinterpreted at the per-market
level: each per-method test rejects a method-specific null about a
method-specific functional of the language-price relationship, while
the joint Hotelling test rejects the global null that no functional
has effect in the given market.

Under this design, the low pairwise Spearman correlations (Shen-Baur
$\rho = -0.15$ and Shen-CF $\rho = -0.07$, neither distinguishable
from zero at our panel size) are exactly the discriminant-validity
pattern the design targets; the moderate Baur-CF correlation
($\rho = +0.48$), with both methods exploiting global semantic
geometry, is the corresponding convergent-validity signal.  The
Hotelling-$T^2$ rejection of the joint null in eleven of twelve
markets, while any individual method rejects in only a small
subset, is the canonical signature predicted when several low-power
estimators each identify a different feature of a partially-
identified estimand \citep{manski2003partial, tamer2010partial}: the
joint test inherits power from the union of identifying restrictions
even when each marginal test does not.

The second reading attributes the discordance to model
misspecification. The Shen-Ross channel is essentially a one-
dimensional summary of the listing's language; the Baur PC1 channel
is essentially a one-dimensional summary of a $768$-dimensional
encoder; the counterfactual channel measures an intervention that
the rewriting model may not execute perfectly. The estimands these
channels target may all be the same underlying causal quantity but
the estimators applied to our finite sample produce different
point estimates because of finite-sample noise and method-specific
specification error. Under this reading the cross-method discordance
is evidence that the variance of language-as-treatment estimation
exceeds the bias of any one method, and the right response is to
deploy all three jointly and report the union or intersection of
identified markets as the conservative or aggressive answer.

The third reading is intermediate. The three channels probably
do not all target the same causal quantity, but the overlap is
positive: $3$ of the $12$ markets are identified at the $95\%$ level
by at least two of the three methods, with no market identified by
all three. San Francisco and Dallas are identified by Shen and the
counterfactual; Portland is identified by Baur and the
counterfactual.  Under the pooled-PCA Baur specification of
Section~\ref{sec:baur-12city}, only Portland returns an individual
$95\%$ confidence interval that excludes zero on the Baur axis, so
the multi-method identification overlap on the Baur side concentrates
in that single market; the joint-null Hotelling-$T^2$ test below
takes up the question of cross-method consistency more directly and
yields a much larger rejection set.

\paragraph{Joint-null Hotelling test.} A more powerful summary of
cross-method evidence is to ask whether, for each market, the joint
null  H_0: (\theta_{\mathrm{Shen}}, \theta_{\mathrm{Baur}},
\theta_{\mathrm{CF}}) = (0, 0, 0) can be rejected against the two-
sided alternative.  Treating the three per-market estimates as
asymptotically independent (the three methods use disjoint sources
of variation: a Doc2Vec uniqueness score, a sentence-BERT principal
component, and an LLM-rewriting differential), we compute the
Hotelling-style statistic $T^2_m = z^2_{\mathrm{Shen},m} +
z^2_{\mathrm{Baur},m} + z^2_{\mathrm{CF},m}$ per market and refer it
to a $\chi^2_3$ asymptotic distribution.

The joint null is rejected at conventional $5\%$ significance in
eleven of the twelve markets and survives only in Atlanta
($T^2 = 3.07$, $p = 0.38$). Under Benjamini-Hochberg false-discovery
correction across the twelve markets, the joint null is rejected in
eleven of twelve at $q < 0.02$. The Hotelling rejection in eleven
markets despite weak pairwise rank correlations between the
individual methods is consistent with each method contributing
independent low-power evidence of a non-zero text-on-price signal,
even when no single method attains significance on that market alone.
Atlanta is the only market where all three methods return small
estimates whose joint null cannot be rejected, and we interpret
Atlanta as a clean null for the text-on-price relationship in the
post-2020 listing corpus.

We favor the third reading for the JBES revision because it is the
most parsimonious of the three and because it preserves the
substantive finding that text-on-price effects are real but
heterogeneous across metropolitan markets, while leaving open the
question of which method is preferred when an analyst can pick only
one.

\subsubsection{Joint identification of the latent text-price effect}
\label{sec:joint-identification}

The cross-method pattern reported above admits a constructive
interpretation as well as the discriminant-validity reading.  We
treat the three per-market point estimates as noisy indicators of a
single latent market-level text-price effect $\theta_m$ and estimate
$\theta_m$ via an anchored Bayesian one-factor model in the tradition
of \citet{joreskog1969general} and \citet{bartholomew2011latent}.
Let $z_{km}$ denote the MAD-standardized point estimate from method
$k$ in market $m$ with reported standard error $s_{km}$. The
measurement equation is
\begin{equation}
z_{km} = \lambda_k \, \theta_m + \varepsilon_{km}, \qquad
\varepsilon_{km} \sim \mathcal{N}\bigl(0, \sigma_k^2 + s_{km}^2\bigr),
\label{eq:cfa-measurement}
\end{equation}
with the structural equation
$\theta_m \sim \mathcal{N}(\mu + \beta' X_m, \tau^2)$ in the
moderators $X_m$ of Section~\ref{sec:meta-regression-12}.  We impose
the identification restriction
$\lambda_{\mathrm{Baur}} = \lambda_{\mathrm{CF}} = +1$: the Baur
pooled-PCA axis and the LLM counterfactual intervention both exploit
global semantic geometry, so the model anchors their shared loading
sign and lets the Shen loading $\lambda_{\mathrm{Shen}}$ be
estimated freely as a discriminant-validity test.  The posterior of
$\theta_m$ given the three-vector
$(z_{\mathrm{Shen},m}, z_{\mathrm{Baur},m}, z_{\mathrm{CF},m})$ is
the precision-weighted closed-form normal whose mean is the joint
estimator we report; the posterior SE inherits both the
method-specific sampling SEs $s_{km}$ and the residual
method-specific variances $\sigma_k^2$.

The fitted loading on the discriminant axis is
$\widehat{\lambda}_{\mathrm{Shen}} = -0.44$, with residual standard
deviations $\widehat{\sigma}_{\mathrm{Shen}} = 0.03$,
$\widehat{\sigma}_{\mathrm{Baur}} = 3.98$, and
$\widehat{\sigma}_{\mathrm{CF}} = 4.42$ on the standardized scale.
The negative Shen loading is consistent with the substantive reading
that within-market lexical distinctiveness moves the joint latent
factor in the opposite direction from the utilitarian-axis discount
that Baur and the counterfactual co-identify.  Despite weak pairwise
rank correlations between the individual methods, the posterior of
$\theta_m$ excludes zero in all twelve of twelve markets at the
$95\%$ credible-interval level, strengthening the Hotelling-$T^2$
result of eleven-of-twelve rejection by adding Atlanta to the
identified set on the basis of cross-method information borrowing.
The combination is formally equivalent to a Stouffer-Lipt\'ak
weighted $Z$-pool of the three method-level test statistics
\citep{liptak1958combination}; we report both summaries for
transparency.  The interpretation is the one
\citet{heckman2008econometric} anticipates: each method identifies a
method-specific functional of the language-price relationship, and
the joint model recovers the common scalar factor on which all three
functionals load, with method-specific loadings encoding the
discriminant-versus-convergent signal in a single estimator.
